% The structure-elucidation problem as a set map: f : M -> S is neither injective
% (m1,m2,m3 land inside one epsilon-ball, so no measurement separates them) nor
% surjective (one spectrum has no preimage), hence f^{-1} is set-valued.
% Requires: shared/notation.tex (fig styles, amber), tikz libs arrows.meta, calc.
\begin{tikzpicture}[x=1cm,y=1cm,fig,>={Latex[length=2.8mm,width=2.0mm]},
  % --- glyphs, declared picture-locally so re-\input stays safe -------------
  molecule/.pic={
    \foreach \a/\b in {30/90,90/150,150/210,210/270,270/330,330/30}
      \draw[fig heavy] ({0.26*cos(\a)},{0.26*sin(\a)}) -- ({0.26*cos(\b)},{0.26*sin(\b)});
    \foreach \ang in {30,90,150,210,270,330}
      \filldraw ({0.26*cos(\ang)},{0.26*sin(\ang)}) circle (1.5pt);
  },
  % #1 = four bar heights, so near-identical spectra can differ just slightly
  spectrum/.pic={
    \draw[fig thin] (-0.32,-0.22) -- (0.32,-0.22);
    \foreach \h [count=\i from 0] in {#1}
      \draw[fig line] ({-0.24+0.16*\i},-0.22) -- ({-0.24+0.16*\i},{-0.22+\h});
  },
  molnode/.style={circle,inner sep=0pt,minimum size=0.62cm},
  specnode/.style={rectangle,inner sep=0pt,minimum width=0.68cm,minimum height=0.62cm},
  glyphlbl/.style={font=\scriptsize,black!60},
  inv/.style={BrickRed!55,fig line,dashed,dash pattern=on 2.2pt off 1.8pt,
              -{Latex[length=2.2mm,width=1.5mm]}},
]

  % ===== the two spaces =====================================================
  \draw[black!40,fig thin,fill=black!3,rounded corners=12pt] (-1.30,-5.30) rectangle (1.30,0.20);
  \draw[black!40,fig thin,fill=black!3,rounded corners=12pt] ( 6.10,-5.30) rectangle (9.80,0.20);

  \node[panel title,anchor=south] at (0,0.32)
       {$\mathcal{M}$ {\normalfont\scriptsize\color{black!60}molecular structures}};
  \node[panel title,anchor=south] at (7.95,0.32)
       {$\mathcal{S}$ {\normalfont\scriptsize\color{black!60}mass spectra}};

  % ===== the epsilon-area, inside S =========================================
  \fill[amber!12] (7.65,-1.25) circle (1.38);
  \draw[amber!80!black,fig line,dashed] (7.65,-1.25) circle (1.38);
  \draw[amber!70!black,fig hair] (7.65,-1.25) -- ++(20:1.38);
  \node[font=\scriptsize,amber!55!black] at ($(7.65,-1.25)+(20:0.70)+(0,0.20)$) {$\varepsilon$};

  % ===== molecules ==========================================================
  \node[molnode] (m1) at (0,-0.50) {};
  \node[molnode] (m2) at (0,-1.25) {};
  \node[molnode] (m3) at (0,-2.10) {};
  \node[molnode] (m4) at (0,-4.10) {};

  \foreach \m in {m1,m3,m4} \pic[Cerulean!70] at (\m) {molecule};
  \pic[RoyalBlue!80!black] at (m2) {molecule};

  \node[glyphlbl,anchor=south west]                       at (0.26,-0.30) {$m_1$};
  \node[font=\scriptsize,RoyalBlue!70!black,anchor=south west] at (0.26,-1.05) {$m_2$};
  \node[glyphlbl,anchor=south west]                       at (0.26,-1.90) {$m_3$};
  \node[glyphlbl,anchor=south west]                       at (0.26,-3.90) {$m_4$};

  % ===== spectra ============================================================
  % three near-identical profiles, all inside the epsilon-ball
  \node[specnode] (s1) at (7.35,-0.40) {};
  \node[specnode] (s2) at (8.35,-1.75) {};
  \node[specnode] (s3) at (7.30,-2.20) {};
  \pic[amber!85!black] at (s1) {spectrum={0.30,0.52,0.22,0.40}};
  \pic[amber!85!black] at (s2) {spectrum={0.32,0.50,0.24,0.38}};
  \pic[amber!85!black] at (s3) {spectrum={0.28,0.52,0.20,0.42}};

  % the measured spectrum: the point the ball is drawn around
  \node[circle,inner sep=0pt,minimum size=0.20cm] (sdot) at (7.65,-1.25) {};
  \fill[amber!85!black] (sdot) circle (2.8pt);
  \node[font=\scriptsize,amber!55!black,anchor=south] at (7.60,-1.14) {$s$};

  % a well-separated case, and one spectrum nothing maps to
  \node[specnode] (s4) at (7.35,-4.10) {};
  \pic[amber!85!black] at (s4) {spectrum={0.52,0.18,0.44,0.26}};
  \node[specnode] (s5) at (8.60,-4.10) {};
  \pic[black!35] at (s5) {spectrum={0.20,0.30,0.52,0.46}};
  \node[font=\scriptsize,black!55,anchor=north,align=center] at (8.60,-4.52)
       {no molecule\\produces this};

  % ===== the forward map ====================================================
  \draw[black!65,fig line,->] (m1) -- (s1);
  \draw[RoyalBlue!70!black,fig medium,->] (m2) -- (s2);
  \draw[black!65,fig line,->] (m3) -- (s3);
  \draw[black!65,fig line,->] (m4) -- (s4);

  \node[anchor=south,align=center] at (3.70,0.88)
       {{\small\color{RoyalBlue!70!black}$f : \mathcal{M} \to \mathcal{S}$}\\[1pt]
        {\scriptsize\color{black!55}well-posed: one structure, one spectrum}};

  % what the epsilon-ball means
  \node[font=\scriptsize,amber!45!black,anchor=north,align=center] at (7.65,-2.85)
       {$d(f(m),s)<\varepsilon$\\[1pt] indistinguishable};

  % ===== the inverse: one spectrum, a set of preimages ======================
  \draw[inv] (sdot) to[out=166,in=-10] (m1);
  \draw[inv] (sdot) to[out=180,in=  8] (m2);
  \draw[inv] (sdot) to[out=195,in= -2] (m3);

  \node[font=\scriptsize,BrickRed!60!black,align=center] at (3.70,-3.05)
       {$f^{-1}(s)$\\[1pt] a whole set of candidates};

\end{tikzpicture}
